\documentclass[11pt]{article}

\usepackage{fullpage}

\usepackage{xspace}

\newcommand{\proptitle}[0]{SaTC 2.0: RES: Relational Coverage-Guided Differential Fuzzing for Critical Software}

\pagenumbering{gobble}

\begin{document}

\begin{center}
{\large\sf\textbf{\proptitle}}
\end{center}

\subsection*{Overview}

Semantic inconsistencies between implementations of the same functionality frequently lead to exploitable vulnerabilities and interoperability failures. Differential testing detects such inconsistencies by executing paired implementations on identical inputs and flagging divergent behavior. In practice, modern fuzzers such as AFL++ support this workflow, at best, through a harness that invokes both programs sequentially and compares outputs; coverage from both executions is collapsed into a single linear path, discarding relational structure that could guide exploration toward semantic boundary regions.
 This project advances the hypothesis that differential fuzzing should be treated as first class problem.  Rather than relying on a standard sequential harness and global coverage bitmap, the proposed research develops fuzzing algorithms that explicitly preserve and exploit asymmetries between paired executions. The work will: (1) design relational coverage representations that distinguish novelty in each implementation and in their joint behavior; (2) develop asymmetry-aware mutation and scheduling strategies that prioritize inputs disproportionately influencing one implementation; and (3) provide lightweight support for constructing relational oracles, including input/output normalization layers, for both users desiring a high degree of automation and users requiring fine-grained control via a DSL for differential harnesses.
Evaluation will compare relationally guided fuzzing against the state-of-practice sequential-harness baseline on representative security-relevant targets, including subjects of previous specialized differential fuzzing work and standard benchmarks differentiated via mutation/version selection. 

Keywords: secure software development; differential fuzzing; harness and oracle construction

\subsection*{Intellectual Merit}

The project reframes differential fuzzing as a relational search problem and introduces principled mechanisms for representing and exploiting relational coverage structure. By modeling asymmetric coverage, input sensitivity, and relational outcome classes explicitly, the proposed techniques expose guidance signals that are invisible under conventional sequential coverage aggregation. The research will provide both algorithmic designs and systematic empirical evidence clarifying when relational coverage guidance yields measurable advantages over existing practice. This work contributes new foundations at the intersection of fuzz testing, systems security, and relational program analysis, as well as contributing to the more complex and subtle problem of effectively defining or extracting imperfect differential oracle relations.

\subsection*{Broader Impacts}

Inconsistencies between implementations routinely undermine the security and reliability of widely deployed software infrastructure. By strengthening differential fuzzing against realistic baselines, this project improves practical methods for detecting semantic flaws in security-critical systems. The resulting techniques will be released as open-source extensions to established fuzzing frameworks, facilitating adoption by industry and the research community. The project will train graduate and undergraduate students in secure software analysis and will integrate relational fuzzing concepts into security and systems curricula, contributing to workforce development in software security.

\end{document}
